\addcontentsline{toc}{section}{СПИСОК ИСПОЛЬЗОВАННЫХ ИСТОЧНИКОВ}

\begin{thebibliography}{9}
    \bibitem{zakon1} Российская Федерация. Законы. О защите прав потребителей: Закон РФ от 7 фервраля 1992г N 2300-1 -- Текст : электронный // КонсультантПлюс: справочно-правовая система. URL:https://www.consultant.ru/document/cons\_doc\_LAW\_305/ (дата обращения: 04.05.2026).
    \bibitem{zakon2} Российская Федерация. Законы. Глава 34 «Аренда». Гражданский кодекс РФ -- Текст : электронный // КонсультантПлюс: справочно-правовая система. — URL: https://www.consultant.ru/document/cons\_doc\_LAW\_9027
    /cac65a05697011caaa0eff7d476d867137f1f92a/ (дата обращения: 04.05.2026).
    \bibitem{zakon3} Российская Федерация. Законы. О проведении эксперимента по установлению специального налогового режима "Налог на профессиональный доход": Федеральный закон от 27.11.2018 № 422-ФЗ -- Текст : электронный // КонсультантПлюс: справочно-правовая система. URL:https://www.consultant.ru/document/cons\_doc\_LAW\_311977/ (дата обращения: 04.05.2026).
    \bibitem{zakon4} Российская Федерация. Законы. О персональных данных: Федеральный закон от 27.07.2006 № 152-ФЗ -- Текст : электронный // КонсультантПлюс: справочно-правовая система. URL:https://www.consultant.ru/document/cons\_doc\_LAW\_61801/ (дата обращения: 05.05.2026).
    \bibitem{gost90} ГОСТ 19.701-90. Единая система программной документации. Схемы алгоритмов, программ, данных и систем.– Москва : Стандартинформ, 2010. – Текст~: непосредственный.
    \bibitem{gost602} ГОСТ 34.602-2020. Информационные технологии. Комплекс стандартов на автоматизированные системы. Техническое задание на создание автоматизированной системы.– Москва : Стандартинформ, 2020. – Текст~: непосредственный.
    \bibitem{gost201} ГОСТ 34.201-2020. Информационные технологии. Комплекс стандартов на автоматизированные системы. Виды,комплектность и обозначение документов при создании автоматизированных систем.– Москва : Стандарт информ, 2020. – Текст~: непосредственный.
    \bibitem{django} Дронов, В. А. Django 4. Практика создания веб-сайтов на Python. — СПб.: БХВ-Петербург, 2023. — 800 с. -- ISBN 978-5-9775-1774-4. -- Текст~: непосредственный.
    \bibitem{django} Агафонов, А. А. Основы технологий баз данных: учебное пособие / А.А. Агафонов, А.М. Белов. – Самара: Издательство Самарского университета, 2023. – 304 с. -- ISBN 978-5-7883-1915-5. -- Текст~: непосредственный.
    \bibitem{uml2} Буч, Г. Язык UML. Руководство пользователя / Г. Буч, Д. Рамбо, А. Якобсон.– Москва : ДМК Пресс, 2006.– 496 с.– ISBN 5-94074-334-X. -- Текст~: непосредственный.
    \bibitem{uml1} Новиков, Ф. А. Учебно-методическое пособие по дисциплине «Анализ и проектирование на UML»:учебно-методическое пособие/Ф.А.Новиков. — Санкт-Петербург : НИУ ИТМО, 2007. — 286 с. -— Текст~: электронный.
    \bibitem{apibook} Арно, Л. Проектирование веб-API / Пер. с англ. Д. А. Беликова.– М.: ДМК Пресс, 2020.– 440 с. -- ISBN 978-5-97060-861-6. -- Текст~: непосредственный.
    \bibitem{algoritm} Бхаргава, А. Грокаем алгоритмы. Иллюстрированное пособие для программистов и любопытствующих. - СПб.: Питер, 2017. - 288 с. -- ISBN 978-5-496-02541-6. -- Текст~: непосредственный.
    \bibitem{dockerbook} Милл, И. Docker на практике / пер. с англ. Д.А. Беликов. – М.: ДМК  Пресс, 2020. – 516 с. -- ISBN 978-5-97060-772-5. -- Текст~: непосредственный.
    \bibitem{pythonbookneed} Прохоренок, Н. А. Python 3. Самое необходимое / Н. А. Прохоренок, В. А. Дронов. — 2-е изд., перераб. и доп. — СПб.: БХВ-Петербург, 2019. — 608 с. -- ISBN 978-5-9775-3994-4. -- Текст~: непосредственный. 
    \bibitem{pythonknig} Мэтиз, Э. Изучаем Python: программирование игр, визуализация данных, веб-приложения. 3-е изд. — СПб.: Питер, 2020 — 512 с. -- ISBN 978-5-4461-1528-0. -- Текст~: непосредственный.
    \bibitem{goodbook} Чаплыгин, А.А. Основы дипломного проектирования: учебное пособие для студентов направления подготовки 09.03.04 Программная инженерия, направленность (профиль) «Разработка программноинформационных систем» / А.А. Чаплыгин, В.В. Апальков, А.В. Малышев ; Юго-Зап. гос. ун-т. – Курск, 2022. – 102 с. – Библиогр.: 76–79 с. -- ISBN 978-5-907407-65-7. -- Текст~: непосредственный.
    \bibitem{restbook} Амундсен, М. RESTful Web API. Паттерны и практики. — Астана: «Спринт Бук», 2025. — 464 с. -- ISBN 978-601-08-4867-2. -- Текст~: непосредственный. 
    \bibitem{recursebook} Свейгарт, Э. Рекурсивная книга о рекурсии. — СПб.: Питер, 2023. — 336 с. -- ISBN 978-5-4461-2393-3. - Текст~: непосредственный.
    \bibitem{webapi} Аттарди, Д. Web API. Сборник рецептов: пер. с англ. - Астана: АЛИСТ, 2025. 304 с. -- ISBN 978-601-12-3681-2. - Текст~: непосредственный. 
    \bibitem{apipaint} Константинов, С. С. АРI как искусство: разработка, поддержка, интеграция. - СПб.: БХВ-Петербурr, 2025. - 304 с. -- ISBN 978-5-9775-2084-3. -- Текст~: непосредственный. 
    \bibitem{apivladik} Светлаков В. И. Архитектура бэкенда. APl для надежных корпоративных приложений. СПб.: БХВ-Петербург, 2026. - 256 с. -- ISBN 978-5-9775-2102-4. -- Текст~: непосредственный. 
    \bibitem{pythonsimple} Нисчал Н. Python — это просто. Пошаговое руководство по программированию и анализу данных: Пер. с англ. — СПб.: БХВ-Петербург, 2023. — 416 с. -- ISBN 978-5-9775-6849-4. -- Текст~: непосредственный.
    \bibitem{pythonbook3} Левашов, П. Python с нуля. — СПб.: Питер, 2024. — 448 с. -- ISBN 978-5-4461-2145-8. -- Текст~: непосредственный
    \bibitem{pythonbook4} Афанасьев, А. Ю. Веб-программирование : учебник и практикум для академического бакалавриата / А. Ю. Афанасьев. — 2-е изд., перераб. и доп. — Москва: Юрайт, 2019. — 333 с. —- ISBN 978-5-534-01432-5. -- Текст~: непосредственный.
    \bibitem{oopbook} Кальб, И. Объектно-ориентированное программирование с помощью Python / Ирв Кальб ; [перевод с английского М. А. Райтмана]. — Москва : Эксмо, 2024. — 512 с. —- ISBN 978-5-04-186627-3. -- Текст~: непосредственный.
    \bibitem{receiptbook} Юн,	Ц. Рецепты Python. Коллекция лучших техник программирования. — СПб.: Питер, 2024. — 512 с. -- ISBN 978-5-4461-2156-4. -- Текст~: непосредственный.
    \bibitem{pssql} Моргунов, Е. П. PostgreSQL. Основы языка SQL/ Е. П. Моргунов. —Санкт-Петербург : БХВ-Петербург, 2018. — 336 с. —- ISBN-978-5-9775-3697. -- Текст~: непосредственный.
    \bibitem{ingir} Соммервиль, И.Инженерия программного обеспечения /И. Соммервиль ; пер. с англ. — 10-е изд. — Москва : Вильямс, 2019. — 1072 с. -- ISBN-978-5-907036-78-5. - Текст~: непосредственный.
    \bibitem{mashtabapp} Атчисон, Л. Масштабирование приложений. Выращивание сложных систем. - СПб.: Питер, 2018. - 256 с. -- ISBN 978-5-496-02952-0 -- Текст~: непосредственный.
    \bibitem{github} Дронов В. А. Фишерман Л. В. Git. Практическое руководство. Управление и контроль версий в разработке программного обеспечения - Спб.: Наука и техника, 2021. - 304 с. -- ISBN 978-5-94387-547-2 -- Текст~: непосредственный.
    \bibitem{pythontt} Кольцов, Д. В. Рython: Создаем программы и игры. 2-е издание. - Спб.: Наука и Техника, 2019.-400 с. -- ISBN 978-5-94387-778-0. –- Текст~: непосредственный. 
    \bibitem{viger} Вигерс, К. Разработка требований к программному обеспечению: [практические приёмы сбора требований и управления ими при разработке программных продуктов : пер. с англ.] / К. Вигерс, Д. Битти. – 3-е изд., доп. – Санкт-Петербург : BHV, 2020. – 736 с. -- ISBN 978-5-9775-3348-5. –- Текст~: непосредственный.
    \bibitem{martin} Мартин, Р. Чистая архитектура. Искусство разработки про-граммного обеспечения / Р. Мартин ; пер.  с англ. А. Кисилева. – Санкт-Петербург : Питер, 2018. – 351 с. – ISBN 978-5-4461-0772-8. –- Текст~: непосредственный.
    \bibitem{belik} Белик, А. Г. Проектирование и архитектура программных систем : учебное пособие / А. Г. Белик, В. Н. Цыганенко. – Омск : ОмГТУ, 2016. – 96 с. –- ISBN 978-5-8149-2258-8. –- Текст~: непосредственный.
    \bibitem{nginx} Дерек де Йонге. NGINX. Книга рецептов. / пер. с англ. Д. А. Беликова. – М.: ДМК Пресс, 2020. – 176 с. -- ISBN 978-5-97060-790-9. -- Текст~: непосредственный.
    \bibitem{postgre}	Ткачев, О. А. PostgreSQL: SQL + PL/pgSQL для тех, кто хочет стать профессионалом — СПб.: Издательство Наука и Техника, 2024. — 480 с. -- ISBN 978-5-907592-32-2. -- Текст~: непосредственный.
	\bibitem{pythonbook} Лейси, Н. Python, например. — СПб.: Питер, 2021. — 208 с. — ISBN 978-5-4461-1826-7. -- Текст~: непосредственный.
    \bibitem{sqlbook} Вальдуриес, П. Принципы организации распределенных баз данных / пер. с  англ. А. А. Слинкина. – М.: ДМК Пресс, 2021. – 672 с. -- ISBN 978-5-97060-391-8 — Текст~: непосредственный.
    \bibitem{webbook} Пьюривал, С. Основы разработки веб-приложений. — СПб.: Питер, 2015. — 272 с. -- ISBN 978-5-496-01226-3. -- Текст~: непосредственный. 
    \bibitem{postgresite} PostgreSQL Documentation : официальный сайт. – PostgreSQL, 1996 – URL: https://www.postgresql.org/docs/ (дата обращения: 07.05.2026).

\bibitem{djangosite} Django Documentation : официальный сайт. – Django Software Foundation, 2005 – URL: https://docs.djangoproject.com/en/6.0/ (дата обращения: 04.05.2026).

\bibitem{pythondoc} Python Documentation : официальный сайт. – Python Software Foundation, 2001 – URL: https://www.python.org/doc/ (дата обращения: 11.05.2026).

\bibitem{gitsite} Git Documentation : официальный сайт. – Software Freedom Conservancy, 2005 – URL: https://git-scm.com/docs/git (дата обращения: 11.05.2026).

\bibitem{dockersite} Docker Documentation : официальный сайт. – Docker Inc., 2013 – URL: https://docs.docker.com (дата обращения: 12.05.2026).

\bibitem{vscode} Visual Studio Code Documentation : официальный сайт. – Microsoft, 2015 – URL: https://code.visualstudio.com/docs (дата обращения: 05.05.2026).

\bibitem{nginxsite} Nginx Documentation : официальный сайт. – F5, Inc., 2004 – URL: https://nginx.org/ru/docs/ (дата обращения: 14.05.2026).

\bibitem{gunsite} Gunicorn Documentation : официальный сайт. – Gunicorn, 2010 – URL: https://gunicorn.org/deploy/ (дата обращения: 14.05.2026).

\bibitem{swagsite} Swagger Documentation : официальный сайт. – SmartBear Software, 2011 – URL: https://docs.swagger.io/spec.html (дата обращения: 14.05.2026).

\bibitem{httpsite} HTTP Documentation : официальный сайт. – HTTP Working Group, 1999 – URL: https://httpwg.org/specs/ (дата обращения: 14.05.2026).
\end{thebibliography}
